Je werkt als systeembeheerder bij een klein ICT-bedrijf.

Binnen het bedrijf werken verschillende medewerkers. Medewerkers moeten bestanden kunnen delen met collega's die tot hetzelfde team behoren. Bestanden mogen niet zomaar toegankelijk zijn voor medewerkers uit andere teams.

Je gaat daarom een eenvoudige gebruikers- en rechtenstructuur maken in Linux.

Je gaat werken met:
\begin{itemize}
\item gebruikers;
\item groepen;
\item bestandseigenaars;
\item groepsrechten;
\item lees-, schrijf- en uitvoerrechten;
\item \texttt{chmod};
\item \texttt{chown};
\item \texttt{chgrp};
\item het \textbf{setgid-bit} op een directory.
\end{itemize}

\textbf{Let op:} voer opdrachten waarbij je gebruikers en groepen aanmaakt uit met voldoende rechten, bijvoorbeeld via \texttt{sudo}.

